%=========== ΛΥΣΗ - 10ο ΘΕΜΑ Α ===========
\providecommand{\theoriaref}[3]{%
\begin{center}
\fcolorbox{maincolor!70!black}{maincolor!8}{%
\begin{minipage}{.88\linewidth}
\textcolor{maincolor!80!black}{\kerkissans{\textbf{#1~\ref{#2}}}}\\[1mm]
#3
\end{minipage}}
\end{center}}
\providecommand{\orismosref}[1]{%
\theoriaref{Ορισμός}{#1}{Η απάντηση δίνεται από τον αντίστοιχο ορισμό.}}
\providecommand{\apodeixiref}[1]{%
\theoriaref{Θεώρημα}{#1}{Η ζητούμενη απόδειξη δίνεται στην αντίστοιχη απόδειξη.}}

\begin{thema}{Α}
\begin{erwthma}
\item \apodeixiref{thewrhma:endiameswn_timwn}

\item \orismosref{orismos:deyterh_paragwgos}

\item \orismosref{orismos:arxikh_synarthsh}

\item Οι χαρακτηρισμοί των προτάσεων είναι:
\begin{alist}
\item \textbf{Σωστό}. Αφού η $f$ παίρνει και θετικές και αρνητικές τιμές, υπάρχουν τμήματα με μη μηδενικό εμβαδόν πάνω ή κάτω από τον άξονα. Το ολοκλήρωμα μηδενίζεται λόγω συμψηφισμού, ενώ το γεωμετρικό εμβαδόν είναι θετικό.
\item \textbf{Λάθος}. Από το $f'(\xi)=0$ για κάποιο $\xi\in(a,\beta)$ δεν προκύπτει ότι $f(a)=f(\beta)$. Για παράδειγμα, για $f(x)=x^2$ στο $[-1,2]$ έχουμε $f'(0)=0$, αλλά $f(-1)\neq f(2)$.
\item \textbf{Λάθος}. Η $f(x)=\sqrt{x}$ δεν είναι παραγωγίσιμη στο $0$, αφού
\[ \lim_{x\to0^+}\dfrac{\sqrt{x}-0}{x-0}=\lim_{x\to0^+}\dfrac{1}{\sqrt{x}}=+\infty. \]
\item \textbf{Σωστό}. Το πεδίο ορισμού του αθροίσματος $f+g$ είναι το κοινό πεδίο ορισμού των $f,g$, δηλαδή $D_f\cap D_g$.
\item \textbf{Λάθος}. Ίσα ολοκληρώματα δεν σημαίνουν ίσες συναρτήσεις. Για παράδειγμα, στο $[-1,1]$ έχουμε
\[ \dint_{-1}^{1}x\d x=0=\dint_{-1}^{1}0\d x, \]
αλλά $x\neq0$ για κάθε $x$.
\end{alist}

\item Σωστή απάντηση είναι η \textbf{β}. Μια συνεχής συνάρτηση μπορεί να έχει ακρότατο σε σημείο όπου δεν υπάρχει παράγωγος, όπως η $f(x)=|x|$ στο $0$.
\end{erwthma}
\end{thema}
%=========== ΤΕΛΟΣ ΛΥΣΗΣ - 10ο ΘΕΜΑ Α ===========
